% fig:acfg --- the ACFG as a data structure: a node tree plus sidecar tables.
% Faithful to chapters/06-architecture.tex sec:arch-acfg:
%   nodes = operations (kernel firing + dataflow edges), sequences, repeats
%           (loops w/ back-edge), later joined by sync + transfer nodes.
%   sidecars = auxiliary maps kept BESIDE the tree, keyed by node identity:
%     per-worker loop ranges (partition), per-axis halo widths (halo),
%     reuse-window slots (reuse), grid shape + neighbour pairs (2D partition),
%     initial-marking depth (pipeline). Generic shape; the stencil instance is
%     fig:wt-acfg.
\begin{figure}[tbp]
  \centering
  \begin{tikzpicture}[
      font=\small,
      seq/.style={draw, rounded corners, fill=gray!15, inner sep=3pt, minimum height=6mm},
      op/.style={draw, fill=blue!10, inner sep=3pt, minimum height=6mm},
      rep/.style={draw, fill=orange!14, inner sep=3pt, minimum height=6mm},
      sx/.style={draw, fill=green!12, inner sep=3pt, minimum height=6mm, font=\footnotesize},
      t/.style={-{Latex[length=1.5mm]}},
      back/.style={-{Latex[length=1.5mm]}, densely dashed, gray!55!black},
    ]
    % ---- node tree (generic): Op/Repeat are authored; Sync/Xfer are injected
    %      as PEER nodes of the enclosing Seq (not nested under one another). ----
    \node[seq] (seq) at (0,0) {\texttt{Seq}};
    \node[op]  (op1) at (-3.0,-1.3) {\texttt{Op}};
    \node[rep] (rep) at (-1.0,-1.3) {\texttt{Repeat}};
    \node[sx, fill=green!12]  (sync) at (1.0,-1.3) {\texttt{Sync}};
    \node[sx, fill=green!12] (xfer) at (3.0,-1.3) {\texttt{Xfer}};
    \node[op]  (op2) at (-1.0,-2.6) {\texttt{Op}};
    \draw[t] (seq) -- (op1);
    \draw[t] (seq) -- (rep);
    \draw[t] (seq) -- (sync);
    \draw[t] (seq) -- (xfer);
    \draw[t] (rep) -- (op2);
    \draw[back] (op2.west) to[out=180,in=180,looseness=1.8]
      node[left=1pt,font=\scriptsize,pos=0.5]{back-edge} (rep.west);

    \node[font=\footnotesize\bfseries] at (0,0.9) {node tree};
    \node[font=\scriptsize, align=left, anchor=north west] at (-3.7,-3.3)
      {\textit{node kinds:} \texttt{Op} (kernel firing $+$ read/write dataflow edges),
       \texttt{Seq}, \texttt{Repeat} (loop with back-edge); the green
       \texttt{Sync} (barrier)\\ and \texttt{Xfer} (push/wait) are \emph{injected}
       peer nodes, added by the sync/transfer passes.};

    % ---- sidecar tables ----
    \node[draw, fill=yellow!10, rounded corners, align=left, font=\scriptsize,
          text width=52mm, inner sep=4pt, anchor=north west] (side) at (5.6,0.6)
      {\textbf{Sidecar tables} (kept beside the tree, keyed by \emph{any} node's id):\\[2pt]
       $\bullet$ per-worker loop ranges \hfill (partition)\\
       $\bullet$ per-axis halo widths \hfill (halo)\\
       $\bullet$ reuse-window slots \hfill (reuse)\\
       $\bullet$ grid shape $+$ neighbour pairs \hfill (2D partition)\\
       $\bullet$ initial-marking depth \hfill (pipeline)};
    \draw[densely dotted, gray!55!black] (xfer.east) to[out=0,in=180]
      node[above=1pt,font=\scriptsize]{keyed by node id} (side.west);
  \end{tikzpicture}
  \caption[The ACFG data structure]{The application control-flow graph (ACFG):
    a tree whose nodes are operations (a kernel firing with its read/write
    dataflow edges), sequences, and repeats (loops with a back-edge), later
    joined by the injected synchronisation and transfer nodes (green, added as
    peer nodes of the enclosing sequence by the sync and transfer passes). Rather
    than rewriting the tree, each middle-end pass populates \emph{sidecar}
    tables ---
    auxiliary maps kept beside the tree and keyed by node identity --- so a pass
    adds precisely the information it derives (per-worker ranges, halo widths,
    reuse slots, grid/neighbour data, pipeline initial-marking depth) without
    disturbing what earlier passes established. The concrete instance for the
    stencil is shown in \cref{fig:wt-acfg}.}
  \label{fig:acfg}
\end{figure}
